\begin{abstract}
Many agent tasks require long action sequences whose effects compound and whose result can be
checked only at the end. We introduce \textbf{CP2Seq}, a benchmark for long-horizon sequence planning
with exact verification, built on flat origami. Each of 600 procedurally generated samples gives a
crease pattern and a target folded state, including layer order; the model must find a sequence of
simple folds from an unfolded square. An engine executes or rejects each fold with a named reason,
and a sequence is accepted only if its final stack matches the target layer by layer, including
which region of the sheet each layer comes from. We make no claim about physical folding. We
evaluate GPT-5.6 and GPT-6 Luna under increasing tool assistance and Claude Sonnet 5.5 on 20
instances, counting exact matches with one episode per instance unless stated. The models solve
some short sequences and none of the 20 medium or hard ablation instances (11 or more folds). To
tell whether this reflects the instances or the models, we run random legal folding over the same
actions and budget: it solves 57\% of easy and 96\% of some-layers episodes, while no model
condition solves more than 5 and 8 of the 10 instances. The targets are reachable; the models'
choices are less effective than chance. Their episodes show why. They rarely reach the target and
fail to stop (2 extra solves among 143 logged episodes), they reach a median tenth to a quarter of
the states random folding reaches, and they restart about as often when a repair within three
folds exists as when none does. Support raises totals without removing this pattern: higher
reasoning effort lifts GPT-6 Luna's some-layers solves from 6/30 to 18/30, and legal-fold lists
and target feedback each help. Because the models are given an exact numeric description of every layer
(coordinates, orientation, stacking order and sheet region), we also test whether rendered images
add anything; removing them does not lower solve rate, so CP2Seq measures planning over that
description rather than visual perception. Single runs vary
substantially, and depth and layer count are confounded in the deepest group; we report both.
\end{abstract}
